\SourceLeaf{68}
The proposed task is thus complete, insofar as it could be carried out by means of series.

The distinguished \textsc{Kl\"ugel} presented a related investigation (on the potential function

\[
(1-2x\cos\varphi+xx)^{-m}
\]

\noindent in Commentat. Societatis huius Tom. XII). The form he considered
differs from the one examined in this short treatise under \textsc{Euler}'s guidance,
though it would be easy to deduce either from the other. Lest I seem to have
merely repeated work already done, let me say that I undertook the whole investigation
to show by an example that treating infinite series through their general terms poses no difficulty
and often leads to general laws more easily than special devices do.
In the main part of this short treatise I therefore used a different method
from the one followed by the distinguished \textsc{Euler} and \textsc{Kl\"ugel}. In the other parts I tried to reduce everything to general expressions
and to give proofs wherever possible.

\textls[140]{Corollary.}

11. It is well known that higher analysis presents series, finite or infinite, whose terms
at first sight appear to follow different laws. This is especially true of series in
which the terms with even indices are formed differently from those with odd indices.
A separate law is then usually given for each kind. Besides their prolixity, such
definitions contain something contrary to the nature and power of analysis. Its excellence
consists precisely in encompassing under a single general expression even very different
values of one quantity, when these occur, and showing how to derive them all from that expression.
Thus, when several values of a single quantity furnish a solution to a geometric problem,
they are all contained in the equation by which we express the relation of the given quantities
to the unknown. Something similar can be demanded of expressions that give the
general terms of series. If the even and
odd terms differ, the formula for the general term, to be complete, must therefore
give one set of values when the term's unspecified index is even and another when it is
odd. I have already explained in the introduction to this short treatise how this can
be accomplished. Indeed, $(-1)^m$ with $m$ even is $=+1$ and with it odd is $=-1$. Thus
